% Part: first-order-logic
% Chapter: models-theories
% Section: set-theory

\documentclass[../../../include/open-logic-section]{subfiles}

\begin{document}

\olfileid[hi]{fol}{mat}{set}

\olsection{समुच्चयों का सिद्धांत}

लगभग समूचे गणित को समुच्चयों के सिद्धांत में विकसित किया जा सकता है। इस
सिद्धांत में गणित विकसित करने के लिए कई बातें आवश्यक हैं। पहली आवश्यकता
संबंध~$\in$ के अभिगृहीतों का समुच्चय है। कई भिन्न अभिगृहीत-तंत्र विकसित किए
गए हैं, जिनमें~$\in$ के गुणधर्म कभी-कभी परस्पर विरोधी भी हैं। चयन के
अभिगृहीत सहित जर्मेलो--फ्रेंकेल समुच्चय सिद्धांत का अभिगृहीत-तंत्र
$\Log{ZFC}$ विशेष स्थान रखता है। यह अब तक सबसे अधिक प्रयुक्त और अध्ययन किया
गया तंत्र है, क्योंकि इसके अभिगृहीत उन लगभग सभी बातों को सिद्ध करने के लिए
पर्याप्त निकलते हैं जिन्हें गणितज्ञ सिद्ध कर सकने की अपेक्षा रखते हैं। पर यह
स्थापित करने से पहले यह स्पष्ट करना आवश्यक है कि गणितज्ञ जिन बातों को व्यक्त
करना चाहते हैं, उन्हें हम \emph{व्यक्त} भी कैसे कर सकते हैं। आरंभ में ही
भाषा में कोई !!{constant} या !!{function} नहीं है। इसलिए पहली दृष्टि में यह
भी स्पष्ट नहीं कि हम विशिष्ट समुच्चयों (जैसे~$\emptyset$ या~$\Nat$),
समुच्चयों पर संक्रियाओं (जैसे~$X \cup Y$ और~$\Pow{X}$), और उनसे भी आगे
संबंधों तथा फलनों जैसी उन रचनाओं की चर्चा कैसे करेंगे जिनमें समुच्चयों के
अतिरिक्त अन्य प्रकार की वस्तुएँ आती हैं।

पहली बात, ``का अवयव है'' ही हमारी रुचि का एकमात्र संबंध नहीं; ``का उपसमुच्चय
है'' लगभग उतना ही महत्त्वपूर्ण लगता है। किंतु ``का उपसमुच्चय है'' को ``का
अवयव है'' के पदों में \emph{परिभाषित} किया जा सकता है। इसके लिए हमें समुच्चय
सिद्धांत की भाषा में ऐसा !!a{formula}~$!A(x, y)$ खोजना होगा जिसे समुच्चयों का
युग्म~$\tuple{X, Y}$ तभी और केवल तभी संतुष्ट करे जब $X \subseteq Y$। पर~$X$,
~$Y$ का उपसमुच्चय ठीक तभी है जब~$X$ के सभी !!{element}s,~$Y$ के भी
!!{element}s हों। अतः हम~$\subseteq$ को इस सूत्र द्वारा परिभाषित कर सकते हैं:
\[
\lforall[z][(z \in x \lif z \in y)]
\]
अब जब भी किसी सूत्र में संबंध~$\subseteq$ का उपयोग करना हो, हम उसके स्थान पर
उस सूत्र का उपयोग कर सकते हैं (जिसमें~$x$ और~$y$ को उपयुक्त रूप से बदलें और
आवश्यक होने पर आबद्ध चर~$z$ का नाम बदलें)। उदाहरणतः समुच्चयों की
विस्तारात्मकता का अर्थ है कि यदि किन्हीं समुच्चयों~$x$ और~$y$ में से प्रत्येक
दूसरे में समाहित हो, तो~$x$ और~$y$ एक ही समुच्चय होने चाहिए। इसे
$\lforall[x][\lforall[y][((x \subseteq y \land y
    \subseteq x) \lif x = y)]]$ द्वारा, अथवा~$\subseteq$ के स्थान पर ऊपर की परिभाषा रखने पर निम्न
सूत्र द्वारा व्यक्त किया जा सकता है:
\[
\lforall[x][\lforall[y][((\lforall[z][(z \in x \lif z \in y)] \land
    \lforall[z][(z \in y \lif z \in x)]) \lif x = y)]].
\]
वास्तव में यह~$\Log{ZFC}$ के अभिगृहीतों में से एक, ``विस्तारात्मकता का
अभिगृहीत'', है।

$\emptyset$ के लिए कोई !!{constant} नहीं, किंतु ``$x$ रिक्त है'' को
$\lnot \lexists[y][y \in x]$ द्वारा व्यक्त कर सकते हैं। तब ``$\emptyset$
विद्यमान है'' !!{sentence}~$\lexists[x][\lnot \lexists[y][y \in
    x]]$ बन
जाता है। यह~$\Log{ZFC}$ का एक और अभिगृहीत है। (ध्यान दें कि विस्तारात्मकता
के अभिगृहीत से यह निष्पन्न होता है कि केवल एक ही रिक्त समुच्चय है।) समुच्चय
सिद्धांत की भाषा में जब भी~$\emptyset$ की चर्चा करनी हो, हम उसे ``ऐसा समुच्चय
है जो रिक्त है और~\dots'' के रूप में लिखेंगे। उदाहरणतः यह व्यक्त करने के लिए
कि~$\emptyset$ हर समुच्चय का उपसमुच्चय है, हम लिख सकते हैं:
\[
\lexists[x][(\lnot\lexists[y][y \in x] \land \lforall[z][x \subseteq
    z])]
\]
जहाँ निस्संदेह~$x \subseteq z$ को भी उसकी परिभाषा से बदलना होगा।

$X \cup Y$ और~$\Pow{X}$ जैसी समुच्चयों पर संक्रियाओं की चर्चा करने के लिए
हमें इसी प्रकार की युक्ति लगानी होगी। समुच्चय सिद्धांत की भाषा में कोई
फलन-चिह्न नहीं है, किंतु हम फलनात्मक संबंधों~$X
\cup Y = Z$ और~$\Pow{X} = Y$
को इस प्रकार व्यक्त कर सकते हैं:
\begin{align*}
& \lforall[u][((u \in x \lor u \in y) \liff u \in z)]\\
& \lforall[u][(u \subseteq x \liff u \in y )]
\end{align*}
क्योंकि~$X \cup Y$ के !!{element}s ठीक वे समुच्चय हैं जो या तो~$X$ के
!!{element}s या~$Y$ के !!{element}s हैं, और~$\Pow{X}$ के !!{element}s ठीक
~$X$ के उपसमुच्चय हैं।
किंतु इससे हम~$x \cup y$ या~$\Pow{x}$ को पद मानकर उपयोग नहीं कर सकते: हम
केवल उन पूरे !!{formula}s का उपयोग कर सकते हैं जो संबंधों~$X \cup Y = Z$ और
$\Pow{X} = Y$ को परिभाषित करते हैं। वास्तव में अभी हम यह नहीं जानते कि ये
संबंध कभी संतुष्ट होते भी हैं, अर्थात् यह नहीं जानते कि सम्मिलन और घात
समुच्चय सदैव विद्यमान हैं। उदाहरणतः !!{sentence}
$\lforall[x][\lexists[y][\Pow{x} = y]]$,~$\Log{ZFC}$ का एक और अभिगृहीत
(घात समुच्चय का अभिगृहीत) है।

अब क्रमित युग्मों या फलनों की चर्चा का क्या? यहाँ हमें समझाना होगा कि क्रमित
युग्मों और फलनों को विशेष प्रकार के समुच्चय कैसे माना जा सकता है। क्रमित
युग्म~$\tuple{x, y}$ को परिभाषित करने की एक विधि उसे समुच्चय
$\{\{x\}, \{x, y\}\}$ मानना है। किंतु पहले की तरह हम इस समुच्चय को नाम देने
वाला कोई !!a{function} प्रस्तुत नहीं कर सकते; केवल संबंध~$\tuple{x, y} = z$,
अर्थात्~$\{\{x\}, \{x, y\}\} = z$, परिभाषित कर सकते हैं:
\[
\lforall[u][(u \in z \liff (\lforall[v][(v \in u \liff v = x)] \lor
  \lforall[v][(v \in u \liff (v = x \lor v = y))]))]
\]
यह कहता है कि !!{element}s~$u$ ठीक वे समुच्चय हैं जो~$z$ के अवयव हैं और जिनका एकमात्र
!!{element}~$x$ है, अथवा जिनके एकमात्र !!{element}s~$x$ और~$y$ हैं (दूसरे
शब्दों में, वे समुच्चय या तो~$\{x\}$ हैं या~$\{x, y\}$)। यह
मिल जाने पर हम आगे की बातें भी कह सकते हैं, जैसे कि~$X \times Y = Z$:
\[
\lforall[z][(z \in Z \liff \lexists[x][\lexists[y][(x \in
        X \land y \in Y \land \tuple{x, y} = z)]])]
\]

फलन~$f \colon X \to Y$ को संबंध~$f(x)
= y$, अर्थात् युग्मों के समुच्चय
$\Setabs{\tuple{x,y}}{f(x) = y}$ के रूप में माना जा सकता है। तब हम कह सकते
हैं कि कोई समुच्चय~$f$,~$X$ से~$Y$ में फलन है यदि (क) वह संबंध~$\subseteq X \times Y$
है, (ख) वह पूर्ण है, अर्थात् प्रत्येक~$x
\in X$ के लिए कोई~$y \in Y$ है जिसके लिए~$\tuple{x, y} \in f$, और (ग) वह फलनात्मक है, अर्थात् जब भी
$\tuple{x, y}, \tuple{x, y'} \in f$, तब~$y = y'$ (क्योंकि फलन के मान
अद्वितीय होने चाहिए)। अतः ``$f$,~$X$ से~$Y$ में फलन है'' को इस प्रकार लिख
सकते हैं:
\begin{align*}
 \lforall[u][(u \in f \lif {}] & \lexists[x][\lexists[y][(x \in X \land y \in
      Y \land \tuple{x, y} = u)]]) \land {}\\
 \lforall[x][(x \in X \lif {}] &
      (\lexists[y][(y \in Y \land \mathrm{maps}(f, x, y))] \land {}\\
& (\lforall[y][\lforall[y'][((\mathrm{maps}(f, x, y) \land
    \mathrm{maps}(f, x, y')) \lif y = y')]]))
\end{align*}
जहाँ~$\mathrm{maps}(f, x, y)$,~$\lexists[v][(v \in f \land
  \tuple{x, y} = v)]$ का संक्षेप है (यह !!{formula} ``$f(x) = y$'' को व्यक्त करता है)।

अब यह व्यक्त करना भी कठिन नहीं कि~$f\colon X \to Y$ !!{injective} है; उदाहरणतः:
\begin{multline*}
 f \colon X \to Y \land \lforall[x][\lforall[x'][((x \in X \land x' \in
     X \land {}]] \\
 \lexists[y][(\mathrm{maps}(f, x, y) \land \mathrm{maps}(f,
        x', y))]) \lif x = x')
\end{multline*}
फलन~$f\colon X \to Y$ तभी और केवल तभी !!{injective} है जब~$f$,~$x, x'
\in X$ को एक ही~$y$ पर प्रतिचित्रित करे तो~$x = x'$। यदि इस सूत्र को
$\mathrm{inj}(f, X, Y)$ से संक्षिप्त करें, तो हम समुच्चय सिद्धांत की भाषा में
कैंटर के प्रमेय जैसा गैर-तुच्छ कथन भी कह सकते हैं:~$\Pow{X}$ से~$X$ में कोई
!!{injective} फलन नहीं है:
\[
\lforall[X][\lforall[Y][(\Pow{X} = Y \lif
    \lnot\lexists[f][\mathrm{inj}(f, Y, X)])]]
\]

यह लग सकता है कि समुच्चय सिद्धांत को एक और अभिगृहीत चाहिए, जो प्रत्येक
परिभाषक गुणधर्म के लिए किसी समुच्चय के अस्तित्व की गारंटी दे। यदि~$!A(x)$
समुच्चय सिद्धांत का ऐसा सूत्र है जिसमें चर~$x$ मुक्त है, तो हम इस !!{sentence}
पर विचार कर सकते हैं:
\[
\lexists[y][\lforall[x][(x \in y \liff !A(x))]].
\]
यह !!{sentence} कहता है कि ऐसा समुच्चय~$y$ है जिसके !!{element}s ठीक वे सभी
$x$ हैं जो~$!A(x)$ को संतुष्ट करते हैं। इस योजना को ``अप्रतिबंधित
समुच्चय-निर्माण सिद्धांत'' कहते हैं। यह बहुत उपयोगी दिखता है; दुर्भाग्य से यह
असंगत है।~$!A(x) \ident \lnot x \in x$ लें; तब समुच्चय-निर्माण सिद्धांत
कहता है:
\[
\lexists[y][\lforall[x][(x \in y \liff x \notin x)]],
\]
अर्थात् यह उन सभी समुच्चयों के समुच्चय का अस्तित्व कहता है जो स्वयं के
!!{element}s नहीं हैं। ऐसा कोई समुच्चय विद्यमान नहीं हो सकता---यही रसेल का
विरोधाभास है। वास्तव में~$\Log{ZFC}$ में इस सिद्धांत का एक प्रतिबंधित और
संगत रूप, पृथक्करण सिद्धांत, होता है:
\[
\lforall[z][\lexists[y][\lforall[x][(x \in y \liff (x \in z \land
      !A(x))]]].
\]

\begin{prob}
निम्न को दिखाने वाली कोई !!a{derivation} देकर सिद्ध कीजिए कि अप्रतिबंधित
समुच्चय-निर्माण सिद्धांत असंगत है:
\[
\lexists[y][\lforall[x][(x \in y \liff x \notin x)]] \Proves \lfalse.
\]
पहले $(A \lif \lnot A) \land (\lnot A \lif A)
\Proves \lfalse$ दिखाना
सहायक हो सकता है।
\end{prob}
\end{document}
